\section{Introducción}\label{sec:intro}
RIPv2 transporta actualizaciones con autenticación opcional. Cuando se habilita, el resumen criptográfico con clave (\emph{keyed message digest}, RFC~2082) es lo único que impide inyectar una actualización falsificada. Un atacante que produzca una actualización con un digest válido puede inyectar / sobrescribir rutas y desviar el tráfico hacia nodos comprometidos.

Este reporte reconstruye el digest (\autoref{sec:reveng}) y, a partir de eso, recupera por fuerza bruta las claves compartidas de seis capturas RIPv2 autenticadas. La \autoref{sec:protocol} describe el protocolo de coordinación TCP entre maestro y trabajadores; la \autoref{sec:partition}, la partición del espacio de claves; la \autoref{sec:death}, la tolerancia a caídas de nodos; la \autoref{sec:bench}, el rendimiento medido; y la \autoref{sec:conclusion} cierra con conclusiones y trabajo diferido.
